\documentclass[12pt]{article}
\usepackage{latexblog}

% ============================================================
% Blog Metadata
% ============================================================
\blogtitle{一个简单的ETL程序}
\blogdate{2020-05-04}
\blogtags{小实现, MySQL, ETL, 闲话编程}
\bloglang{zh}
\blogsummary{}

\begin{document}
\makeblogtitle

这两天闲着没事准备玩一下社工库，网上有很多以前的社工库（这些都可以下载到，但是实际上已经没有什么太大的价值了，因为暴露时间太久，以及相关的网站都已经做了处理，所以别指望能够找到什么有价值的东西），通过社工库可以了解到的一个实际的数据就是，用户的设置的密码大都是什么样子的。我准备看一下搜云社工库，这个库大概4亿多条数据，主要目的是实践一下大量数据的处理。

\section{还原原始数据库}\label{ux8fd8ux539fux539fux59cbux6570ux636eux5e93}

这个库（下载地址请自行搜索）下载完成之后是一个\texttt{1.bak}（SQL SERVER的备份文件）文件，接近30G，网上有详细的教程如何导入，总结起来大致如下:

\begin{itemize}
\tightlist
\item
  安装SQL SERVER 2008 R2（标准版或者更高）
\item
  将1.bak复制到SQL SERVER的实例备份目录，否则会无法导入
\item
  因为还原后的数据库占用空间在130G，所以要保证数据库的存储磁盘空间足够
\end{itemize}

还原完成之后即可得到一个sgk的数据表，里面有423771078条数据。

\section{导入到MySQl}\label{ux5bfcux5165ux5230mysql}

因为最终希望能够在MYSQL中使用数据，所以还得将数据迁移到MySQL中。最初的想法是直接将数据导出到csv，然后再通过MySQL导入csv。但是这个方法尝试了几次之后发现一些问题：

\begin{itemize}
\tightlist
\item
  数据中可能存在一些非可见字符，导致导出后的csv格式混乱
\item
  莫名其妙的问题无法解析
\end{itemize}

索性就直接放弃了这种方式，而是通过手写ETL来完成迁移。于是设计了一个简陋的ETL程序，还挺有意思的。

\subsection{整体架构}\label{ux6574ux4f53ux67b6ux6784}

ETL的主要目的是把数据从一个地方转移迁移到另一个地方，这其中可能进行一些其他的操作比如清洗或者格式转换。整体的思路是这样的：

\begin{itemize}
\tightlist
\item
  有一个source（源数据）和destination（目标数据源）
\item
  ETL执行的时候从源数据查询数据，然后保存到目标数据库中，为了提高效率，每次会查询一批数据，保存的时候也会一批一批保存
\item
  为进一步提高性能，ETL可以设定多个线程同时处理，为了避免冲突，每个线程不会处理同一条数据，是互斥的
\end{itemize}

\subsection{ETL Engine}\label{etl-engine}

\begin{Shaded}
\begin{Highlighting}[]
\KeywordTok{public} \KeywordTok{class}\NormalTok{ Engine }\OperatorTok{\{}
    \KeywordTok{private} \DataTypeTok{static} \DataTypeTok{final} \BuiltInTok{Logger}\NormalTok{ logger }\OperatorTok{=}\NormalTok{ LoggerFactory}\OperatorTok{.}\FunctionTok{getLogger}\OperatorTok{(}\NormalTok{Engine}\OperatorTok{.}\FunctionTok{class}\OperatorTok{);}

    \KeywordTok{private} \DataTypeTok{final} \DataTypeTok{int}\NormalTok{ threads}\OperatorTok{;}
    \KeywordTok{private} \DataTypeTok{final} \DataTypeTok{int}\NormalTok{ totalCount}\OperatorTok{;}
    \KeywordTok{private} \DataTypeTok{final} \DataTypeTok{int}\NormalTok{ batchSize}\OperatorTok{;}

    \KeywordTok{private} \DataTypeTok{final}\NormalTok{ HikariDataSource source}\OperatorTok{;}
    \KeywordTok{private} \DataTypeTok{final}\NormalTok{ HikariDataSource destination}\OperatorTok{;}

    \KeywordTok{private} \DataTypeTok{final} \BuiltInTok{AtomicLong}\NormalTok{ rangeSelector }\OperatorTok{=} \KeywordTok{new} \BuiltInTok{AtomicLong}\OperatorTok{(}\DecValTok{0}\OperatorTok{);}

    \KeywordTok{private} \DataTypeTok{final} \BuiltInTok{List}\OperatorTok{\textless{}}\BuiltInTok{Thread}\OperatorTok{\textgreater{}}\NormalTok{ tasks }\OperatorTok{=} \KeywordTok{new} \BuiltInTok{ArrayList}\OperatorTok{\textless{}\textgreater{}();}

    \KeywordTok{public} \FunctionTok{Engine}\OperatorTok{(}\DataTypeTok{int}\NormalTok{ threads}\OperatorTok{,} \DataTypeTok{int}\NormalTok{ totalCount}\OperatorTok{,} \DataTypeTok{int}\NormalTok{ batchSize}\OperatorTok{)} \OperatorTok{\{}
        \KeywordTok{this}\OperatorTok{.}\FunctionTok{threads} \OperatorTok{=}\NormalTok{ threads}\OperatorTok{;}
        \KeywordTok{this}\OperatorTok{.}\FunctionTok{totalCount} \OperatorTok{=}\NormalTok{ totalCount}\OperatorTok{;}
        \KeywordTok{this}\OperatorTok{.}\FunctionTok{batchSize} \OperatorTok{=}\NormalTok{ batchSize}\OperatorTok{;}
        \KeywordTok{this}\OperatorTok{.}\FunctionTok{source} \OperatorTok{=} \KeywordTok{new} \FunctionTok{HikariDataSource}\OperatorTok{(}\KeywordTok{new} \FunctionTok{HikariConfig}\OperatorTok{(}\StringTok{"src/main/resources/source.properties"}\OperatorTok{));}
        \KeywordTok{this}\OperatorTok{.}\FunctionTok{destination} \OperatorTok{=} \KeywordTok{new} \FunctionTok{HikariDataSource}\OperatorTok{(}\KeywordTok{new} \FunctionTok{HikariConfig}\OperatorTok{(}\StringTok{"src/main/resources/destination.properties"}\OperatorTok{));}
    \OperatorTok{\}}
\end{Highlighting}
\end{Shaded}

这个ETL引擎担负管理任务的角色，有这些参数：

\begin{itemize}
\tightlist
\item
  threads: 可以设定多少个线程同时运行；
\item
  totalCount: 设定一个totalCount作为任务结束的条件（当然也可以在查询不到数据的时候结束)
\item
  batchSize: 每一批的大小
\item
  rangeSelector: 用来控制每个线程处理的数据，这里将直接用id来区分
\end{itemize}

同时直接初始化了两个数据库的连接池，这里选用的是Hikari连接池。

\begin{Shaded}
\begin{Highlighting}[]
\KeywordTok{public} \DataTypeTok{void} \FunctionTok{run}\OperatorTok{()} \KeywordTok{throws} \BuiltInTok{InterruptedException} \OperatorTok{\{}
    \DataTypeTok{final} \BuiltInTok{CountDownLatch}\NormalTok{ countDownLatch }\OperatorTok{=} \KeywordTok{new} \BuiltInTok{CountDownLatch}\OperatorTok{(}\NormalTok{threads}\OperatorTok{);}
    \ControlFlowTok{for} \OperatorTok{(}\DataTypeTok{int}\NormalTok{ i }\OperatorTok{=} \DecValTok{0}\OperatorTok{;}\NormalTok{ i }\OperatorTok{\textless{}}\NormalTok{ threads}\OperatorTok{;}\NormalTok{ i}\OperatorTok{++)} \OperatorTok{\{}
        \BuiltInTok{Thread}\NormalTok{ thread }\OperatorTok{=} \KeywordTok{new} \BuiltInTok{Thread}\OperatorTok{(}\KeywordTok{new} \FunctionTok{Task}\OperatorTok{(}\NormalTok{source}\OperatorTok{,}
\NormalTok{                destination}\OperatorTok{,}
\NormalTok{                totalCount}\OperatorTok{,}
\NormalTok{                rangeSelector}\OperatorTok{,}
\NormalTok{                batchSize}\OperatorTok{,}
\NormalTok{                countDownLatch}\OperatorTok{));}
\NormalTok{        tasks}\OperatorTok{.}\FunctionTok{add}\OperatorTok{(}\NormalTok{thread}\OperatorTok{);}
    \OperatorTok{\}}
\NormalTok{    tasks}\OperatorTok{.}\FunctionTok{forEach}\OperatorTok{(}\BuiltInTok{Thread}\OperatorTok{::}\NormalTok{start}\OperatorTok{);}
\NormalTok{    countDownLatch}\OperatorTok{.}\FunctionTok{countDown}\OperatorTok{();}
\NormalTok{    countDownLatch}\OperatorTok{.}\FunctionTok{await}\OperatorTok{();}
\OperatorTok{\}}
\end{Highlighting}
\end{Shaded}

核心逻辑就是，启动N个线程，然后一直等待到每个线程都结束，即完成了ETL任务。

\subsection{Task}\label{task}

Task对应到每个线程，每个线程处理的任务是一样的，只是处理的数据记录不同。

\begin{Shaded}
\begin{Highlighting}[]
\KeywordTok{public} \DataTypeTok{void} \FunctionTok{run}\OperatorTok{()} \OperatorTok{\{}
\NormalTok{    logger}\OperatorTok{.}\FunctionTok{info}\OperatorTok{(}\StringTok{"ETL task \{\} running..."}\OperatorTok{,} \BuiltInTok{Thread}\OperatorTok{.}\FunctionTok{currentThread}\OperatorTok{().}\FunctionTok{getId}\OperatorTok{());}
    \ControlFlowTok{try} \OperatorTok{(}\BuiltInTok{Connection}\NormalTok{ sourceConn }\OperatorTok{=}\NormalTok{ source}\OperatorTok{.}\FunctionTok{getConnection}\OperatorTok{();}
         \BuiltInTok{Connection}\NormalTok{ destinationConn }\OperatorTok{=}\NormalTok{ destination}\OperatorTok{.}\FunctionTok{getConnection}\OperatorTok{())} \OperatorTok{\{}
        \ControlFlowTok{while} \OperatorTok{(}\KeywordTok{true}\OperatorTok{)} \OperatorTok{\{}
            \DataTypeTok{final} \DataTypeTok{long}\NormalTok{ startId }\OperatorTok{=}\NormalTok{ rangeSelector}\OperatorTok{.}\FunctionTok{getAndAdd}\OperatorTok{(}\NormalTok{batchSize}\OperatorTok{);}
            \DataTypeTok{final} \DataTypeTok{long}\NormalTok{ endId }\OperatorTok{=}\NormalTok{ startId }\OperatorTok{+}\NormalTok{ batchSize}\OperatorTok{;}
            \ControlFlowTok{if} \OperatorTok{(}\NormalTok{startId }\OperatorTok{\textgreater{}}\NormalTok{ totalCount}\OperatorTok{)} \OperatorTok{\{}
\NormalTok{                logger}\OperatorTok{.}\FunctionTok{info}\OperatorTok{(}\StringTok{"Reached end of records:\{\} , task finished"}\OperatorTok{,}\NormalTok{ startId}\OperatorTok{);}
                \ControlFlowTok{break}\OperatorTok{;}
            \OperatorTok{\}}
\NormalTok{            job}\OperatorTok{.}\FunctionTok{doTransfer}\OperatorTok{(}\NormalTok{sourceConn}\OperatorTok{,}\NormalTok{ destinationConn}\OperatorTok{,}\NormalTok{ startId}\OperatorTok{,}\NormalTok{ endId}\OperatorTok{);}
        \OperatorTok{\}}
    \OperatorTok{\}} \ControlFlowTok{catch} \OperatorTok{(}\BuiltInTok{SQLException}\NormalTok{ ex}\OperatorTok{)} \OperatorTok{\{}
\NormalTok{        logger}\OperatorTok{.}\FunctionTok{error}\OperatorTok{(}\StringTok{"Failed to get data source, ex"}\OperatorTok{);}
    \OperatorTok{\}} \ControlFlowTok{finally} \OperatorTok{\{}
\NormalTok{        logger}\OperatorTok{.}\FunctionTok{info}\OperatorTok{(}\StringTok{"ETL task \{\} finished."}\OperatorTok{,} \BuiltInTok{Thread}\OperatorTok{.}\FunctionTok{currentThread}\OperatorTok{().}\FunctionTok{getId}\OperatorTok{());}
\NormalTok{        countDownLatch}\OperatorTok{.}\FunctionTok{countDown}\OperatorTok{();}
    \OperatorTok{\}}

\OperatorTok{\}}
\end{Highlighting}
\end{Shaded}

处理过程也很简单，拿到一个connection之后，根据rangeSelector得到要获取的id范围，然后委派给具体的Job去处理。当超出最大范围的时候，停止该线程。

\subsection{Job}\label{job}

Job对应到具体每条数据该如何传输，会稍微麻烦一点。但本质还是select然后insert，没什么技术含量。

\begin{Shaded}
\begin{Highlighting}[]
\KeywordTok{public} \KeywordTok{class}\NormalTok{ SeTransferJob }\KeywordTok{implements}\NormalTok{ TransferJob }\OperatorTok{\{}
    \KeywordTok{private} \DataTypeTok{static} \DataTypeTok{final} \BuiltInTok{Logger}\NormalTok{ logger }\OperatorTok{=}\NormalTok{ LoggerFactory}\OperatorTok{.}\FunctionTok{getLogger}\OperatorTok{(}\NormalTok{SeTransferJob}\OperatorTok{.}\FunctionTok{class}\OperatorTok{);}
    \KeywordTok{private} \DataTypeTok{static} \DataTypeTok{final} \BuiltInTok{String}\NormalTok{ query }\OperatorTok{=} \StringTok{"select * from sgk where id \textgreater{}? and id \textless{}=?"}\OperatorTok{;}
    \KeywordTok{private} \DataTypeTok{static} \DataTypeTok{final} \BuiltInTok{String}\NormalTok{ insertSqlTemplate }\OperatorTok{=} \StringTok{"insert into se\_record\_\%d(id, user\_name, email, password, salt, source, remark) values (?,?,?,?,?,?,?);"}\OperatorTok{;}

    \AttributeTok{@Override}
    \KeywordTok{public} \DataTypeTok{void} \FunctionTok{doTransfer}\OperatorTok{(}\BuiltInTok{Connection}\NormalTok{ source}\OperatorTok{,} \BuiltInTok{Connection}\NormalTok{ target}\OperatorTok{,} \DataTypeTok{long}\NormalTok{ startId}\OperatorTok{,} \DataTypeTok{long}\NormalTok{ endId}\OperatorTok{)} \OperatorTok{\{}
\NormalTok{        logger}\OperatorTok{.}\FunctionTok{info}\OperatorTok{(}\StringTok{"Transferring records [\{\}, \{\}]"}\OperatorTok{,}\NormalTok{ startId}\OperatorTok{,}\NormalTok{ endId}\OperatorTok{);}

        \ControlFlowTok{try} \OperatorTok{\{}
\NormalTok{            target}\OperatorTok{.}\FunctionTok{setAutoCommit}\OperatorTok{(}\KeywordTok{false}\OperatorTok{);}
        \OperatorTok{\}} \ControlFlowTok{catch} \OperatorTok{(}\BuiltInTok{SQLException}\NormalTok{ throwables}\OperatorTok{)} \OperatorTok{\{}
            \ControlFlowTok{throw} \KeywordTok{new} \BuiltInTok{RuntimeException}\OperatorTok{(}\StringTok{"Cannot open transaction"}\OperatorTok{);}
        \OperatorTok{\}}
        \DataTypeTok{long}\NormalTok{ tableIndex }\OperatorTok{=} \OperatorTok{(}\NormalTok{startId }\OperatorTok{/} \DecValTok{10000000}\OperatorTok{)} \OperatorTok{+} \DecValTok{1}\OperatorTok{;}
        \BuiltInTok{String}\NormalTok{ insertSql }\OperatorTok{=} \BuiltInTok{String}\OperatorTok{.}\FunctionTok{format}\OperatorTok{(}\NormalTok{insertSqlTemplate}\OperatorTok{,}\NormalTok{ tableIndex}\OperatorTok{);}
        \ControlFlowTok{try} \OperatorTok{(}\BuiltInTok{PreparedStatement}\NormalTok{ queryStatement }\OperatorTok{=}\NormalTok{ source}\OperatorTok{.}\FunctionTok{prepareStatement}\OperatorTok{(}\NormalTok{query}\OperatorTok{);}
             \BuiltInTok{PreparedStatement}\NormalTok{ saveStatement }\OperatorTok{=}\NormalTok{ target}\OperatorTok{.}\FunctionTok{prepareStatement}\OperatorTok{(}\NormalTok{insertSql}\OperatorTok{))} \OperatorTok{\{}
\NormalTok{            queryStatement}\OperatorTok{.}\FunctionTok{setLong}\OperatorTok{(}\DecValTok{1}\OperatorTok{,}\NormalTok{ startId}\OperatorTok{);}
\NormalTok{            queryStatement}\OperatorTok{.}\FunctionTok{setLong}\OperatorTok{(}\DecValTok{2}\OperatorTok{,}\NormalTok{ endId}\OperatorTok{);}
            \ControlFlowTok{try} \OperatorTok{(}\BuiltInTok{ResultSet}\NormalTok{ result }\OperatorTok{=}\NormalTok{ queryStatement}\OperatorTok{.}\FunctionTok{executeQuery}\OperatorTok{())} \OperatorTok{\{}
                \ControlFlowTok{while} \OperatorTok{(}\NormalTok{result}\OperatorTok{.}\FunctionTok{next}\OperatorTok{())} \OperatorTok{\{}
\NormalTok{                    Record }\KeywordTok{record} \OperatorTok{=}\NormalTok{ Record}\OperatorTok{.}\FunctionTok{from}\OperatorTok{(}\NormalTok{result}\OperatorTok{);}
                    \ControlFlowTok{if} \OperatorTok{(!}\KeywordTok{record}\OperatorTok{.}\FunctionTok{isValid}\OperatorTok{())} \OperatorTok{\{}
\NormalTok{                        logger}\OperatorTok{.}\FunctionTok{warn}\OperatorTok{(}\StringTok{"Found invaid record:\{\}"}\OperatorTok{,} \KeywordTok{record}\OperatorTok{);}
                    \OperatorTok{\}} \ControlFlowTok{else} \OperatorTok{\{}
                        \KeywordTok{record}\OperatorTok{.}\FunctionTok{attach}\OperatorTok{(}\NormalTok{saveStatement}\OperatorTok{);}
\NormalTok{                        saveStatement}\OperatorTok{.}\FunctionTok{addBatch}\OperatorTok{();}
                    \OperatorTok{\}}
                \OperatorTok{\}}
            \OperatorTok{\}}
\NormalTok{            saveStatement}\OperatorTok{.}\FunctionTok{executeBatch}\OperatorTok{();}
\NormalTok{            target}\OperatorTok{.}\FunctionTok{commit}\OperatorTok{();}
\NormalTok{            logger}\OperatorTok{.}\FunctionTok{info}\OperatorTok{(}\StringTok{"Records [\{\}, \{\}] {-}\textgreater{} se\_record\_\{\} commited"}\OperatorTok{,}\NormalTok{ startId}\OperatorTok{,}\NormalTok{ endId}\OperatorTok{,}\NormalTok{ tableIndex}\OperatorTok{);}
        \OperatorTok{\}} \ControlFlowTok{catch} \OperatorTok{(}\BuiltInTok{SQLException}\NormalTok{ ex}\OperatorTok{)} \OperatorTok{\{}
\NormalTok{            logger}\OperatorTok{.}\FunctionTok{error}\OperatorTok{(}\StringTok{"Unexpected sql error:"}\OperatorTok{,}\NormalTok{ ex}\OperatorTok{);}
            \ControlFlowTok{throw} \KeywordTok{new} \BuiltInTok{RuntimeException}\OperatorTok{(}\NormalTok{ex}\OperatorTok{);}
        \OperatorTok{\}}
    \OperatorTok{\}}
\OperatorTok{\}}
\end{Highlighting}
\end{Shaded}

这里因为数据量太大，做了一个分区的处理，直接按照id来进行分区，每个表控制在千万以下。因此导入完成后，最终会有40多个表，每个表有接近1千万的数据。否则单表过大之后的插入性能会变得很低。

以上就是整个ETL的核心思想，还是有一定的扩展性的，哈哈。


\end{document}
